\section*{5. Eingrenzung des Themas}
Der Fokus der Arbeit liegt auf dem Monitoring-Prozess. Das Dashboard wird erstellt, um die zu überwachenden Metriken übersichtlich darzustellen. Die Projektskizze beinhaltet dabei den Aufbau und Funktionsweise des Dashboards inkl. des Datenflusses bis zur Anzeige.

Es wird analysiert, inwieweit die neue Lösung die Anforderungen erfüllt und die Effizienz beeinflusst. Weiterhin wird es mithilfe eines Interviews validiert. Hierbei wird die Methodik, bis zur praktischen Umsetzung des Interviews und dessen Auswertung behandelt. Die befragte Personengruppe ist auf das Entwicklungsteam beschränkt, welches das Dashboard nutzen wird.

In der Arbeit wird sich ausschließlich auf diesen Anwendungsfall beschränkt. Das heißt, dass weitere Tools und Strategien zur Qualitätssteigerung und Prozessoptimierung nicht weiter untersucht werden. Der Anbieter zur Dashboard-Erstellung wird nach vorheriger Evaluation ausgewählt. Die restlichen Anbieter sowie alternative Verfahren zur Anzeigenoptimierung werden in diesem Kontext nicht weiter betrachtet. Da das Dashboard die bisherige Lösung in ihrer Funktion ablösen soll, ist auch die Bewertung der Inhalte, wie Vollständigkeit und Sinnhaftigkeit der Tests, nicht Teil der Arbeit.